% Generated draft: requires mathematical and source review.
\documentclass[11pt]{article}
\usepackage[margin=25mm]{geometry}
\usepackage{fontspec}
\setmainfont{Latin Modern Roman}
\usepackage{amsmath,amssymb,mathtools,unicode-math}
\setmathfont{Latin Modern Math}
\usepackage{enumitem,needspace,xurl,hyperref}
\setlength{\parindent}{0pt}
\setlength{\parskip}{5pt}
\setlength{\emergencystretch}{3em}
\newcommand{\N}{\mathbb{N}}
\newcommand{\Z}{\mathbb{Z}}
\newcommand{\Q}{\mathbb{Q}}
\newcommand{\R}{\mathbb{R}}
\newcommand{\C}{\mathbb{C}}
\newcommand{\F}{\mathbb{F}}
\newcommand{\A}{\mathbb{A}}
\newcommand{\PP}{\mathbb{P}}
\newcommand{\E}{\mathbb{E}}
\newcommand{\eps}{\varepsilon}
\newcommand{\dd}{\,\mathrm{d}}
\begin{document}
\section*{UnsolvedMath ID: 30005333}
\textbf{Source problem code:} catalogue entry
\subsection{Definitions and mathematical statement}
The supplied record contains the contribution-title fragment attributed to Van Vu, \emph{Roots of random functions: Recent progress and open questions}. It contains no recoverable mathematical question. In particular, it specifies neither a random-function model nor its probability law, domain, codomain, meaning of a root, or proposed assertion and quantifiers. Consequently, no precise mathematical statement or associated definitions can be reconstructed from the supplied evidence.

The supplied statement-verification field labels this entry \texttt{unrecoverable} and states that specific random-root questions are retained separately under source codes \texttt{OWR-12697684-016} and \texttt{OWR-12697684-017}. Their statements are not supplied here, so neither can be substituted for this entry. The supplied \emph{open} label reflects a curation rule for an incomplete fragment; it does not establish the open status of an identifiable mathematical conjecture.
\subsection{Short English statement}
This entry preserves a talk-title fragment attributed to Van Vu about roots of random functions. No mathematical question can be stated precisely from the supplied record; the intended random-function model and target assertion remain unspecified.
\subsection{Sources}
\begin{itemize}
\item Local UnsolvedMath registry, supplied record with permanent ID \(30005333\), exact source code \texttt{catalogue entry}, title \emph{Unspecified Random Function Root Problem}, in the Oberwolfach Reports Open Problems collection. The supplied proposed year is \(2023\); no proposer is recorded.
\item Supplied report citation: \emph{Combinatorics (2023)}, Oberwolfach Reports, pp.~5--89, DOI \texttt{10.4171/owr/2023/1}. \url{https://doi.org/10.4171/owr/2023/1}. The supplied background describes Oberwolfach Reports as documenting mathematical workshops and their open-problem discussions.
\item Supplied contribution-title attribution: Van Vu, \emph{Roots of random functions: Recent progress and open questions}. The record does not supply a complete contribution passage or a specific question.
\item Supplied literature source, labelled ``original OWR report'': \url{https://doi.org/10.4171/owr/2022/57}. This differs from the report DOI in the supplied source citation; the relationship is unresolved.
\item Literature assessment dated exactly \texttt{2026-08-22}: supplied dated curation, not independently verified. It classifies the fragment as open under its stated rule because no complete claim is available to verify, and requests independent verification of time-sensitive status and cited claims. The supplied references do not provide an authoritative citation for a recoverable, specific mathematical assertion in this entry.
\end{itemize}
\end{document}
